\documentclass{report}
\usepackage[utf8]{inputenc}
\usepackage[spanish, es-noshorthands]{babel}
\usepackage{amsmath}
\usepackage{amssymb}
\usepackage{amsthm}
\usepackage{clrscode3e}

\theoremstyle{definition}
\newtheorem{definicion}{Definición}[chapter]
\theoremstyle{plain}
\newtheorem{teorema}{Teorema}[chapter]
\newtheorem{lema}{Lema}[chapter]
\theoremstyle{remark}
\newtheorem{ejemplo}{Ejemplo}[chapter]

\begin{document}

\setcounter{chapter}{2}
\chapter{Vacío de la intersección de una RCG con un lenguaje regular finito}
\label{cap:vacio-interseccion-rcg-regular}

En el capítulo~\ref{cap:rcg-sat} se presentaron las gramáticas de
concatenación de rango y se revisaron antecedentes en los que el SAT se
aborda mediante intersecciones entre lenguajes. En esos antecedentes, la
fórmula booleana se transforma en un lenguaje regular que representa las
interpretaciones verdaderas bajo el supuesto de que cada ocurrencia de
variable es independiente. Después, una gramática impone la consistencia
entre las ocurrencias de una misma variable. La satisfacibilidad se
reduce entonces al problema de decidir si la intersección de ambos
lenguajes es vacía~\cite{fernandezarias2007,aguileralopez2016}.

El paso de las gramáticas libres del contexto y de las gramáticas de
concatenación de rango simples a las RCG no conserva de manera
automática la tratabilidad de ese enfoque. Bertsch y Nederhof prueban
que, para RCG, el vacío de la intersección con un lenguaje
regular arbitrario es indecidible. También prueban que, si el lenguaje
regular es finito y se representa mediante un autómata acíclico, el
problema de decidir si la intersección es no vacía es
NP-completo~\cite{bertsch2001rcg}. En el presente capítulo se formaliza
ese último caso mediante un algoritmo explícito de enumeración de
caminos aceptantes.

La formalización no sustituye el resultado de complejidad de
Bertsch y Nederhof. Su objetivo es precisar qué procedimiento de decisión
se obtiene cuando el autómata es acíclico, demostrar su correctitud y
caracterizar su costo en función del número de caminos aceptantes del
autómata.

El capítulo se organiza de la siguiente manera. En la
sección~\ref{sec:lenguajes-regulares-finitos} se presentan los lenguajes
regulares finitos. En la sección~\ref{sec:automatas-aciclicos} se
formalizan los autómatas acíclicos y sus caminos aceptantes. En la
sección~\ref{sec:problema-vacio-interseccion} se define el problema de
vacío de la intersección. En la sección~\ref{sec:algoritmo-decision} se
presenta el algoritmo de decisión, se demuestra su correctitud y se
analiza su complejidad. En la
sección~\ref{sec:modelacion-sat} se modela el SAT como un caso de ese
problema. Por último, en la sección~\ref{sec:lectura-sat} se lee el
resultado en términos del SAT y se motiva la búsqueda de subclases no
lineales con vacío tratable.

\section{Lenguajes regulares finitos}
\label{sec:lenguajes-regulares-finitos}

El enfoque SAT por intersección trabaja con un lenguaje regular que
representa interpretaciones de una fórmula. Para una fórmula con una
cantidad finita de ocurrencias de variables, ese lenguaje es finito. La presente
sección fija la noción de lenguaje regular finito que será usada en el
resto del capítulo.

\begin{definicion}
Un lenguaje $L$ es \textbf{finito} si existe un entero $k \ge 0$ tal
que $L$ tiene exactamente $k$ cadenas.
\end{definicion}

\begin{definicion}
Un \textbf{lenguaje regular finito} es un lenguaje finito que también
es regular.
\end{definicion}

Una familia de lenguajes finitos puede tener
una representación mediante autómatas de tamaño polinomial y contener,
al mismo tiempo, una cantidad exponencial de cadenas.

Por ejemplo, para cada $n \ge 1$, sea
\[
  L_n = \{0,1\}^n.
\]
El lenguaje $L_n$ es finito y contiene $2^n$ cadenas. No obstante,
existe un autómata finito con $n+1$ estados que reconoce
$L_n$: el autómata avanza de un estado al siguiente al leer $0$ o $1$, y
acepta solo después de leer exactamente $n$ símbolos.

El ejemplo anterior muestra por qué no basta con saber que el lenguaje
regular es finito. La finitud garantiza que una enumeración termina, pero
no garantiza que termine en tiempo polinomial respecto al tamaño de la
representación. Cuando un lenguaje regular es finito, admite una
representación mediante un autómata acíclico, que se define en la
sección siguiente.

\section{Autómatas acíclicos}
\label{sec:automatas-aciclicos}

En el presente capítulo se usa el autómata finito del
capítulo~\ref{cap:preliminares}, con una diferencia: la función de
transición es parcial, no total, es decir, $\delta(q,a)$ puede no estar
definida para algunos pares $(q,a)$. Cuando $\delta(q,a)$ no está
definida, la lectura del símbolo $a$ en el estado $q$ no conduce a
ningún estado.

\begin{definicion}
Sean $p_0,p_1,\ldots,p_m$ estados de $Q$ y $a_1,a_2,\ldots,a_m$
símbolos de $\Sigma$, con $m \ge 0$. Un \textbf{camino} desde $p_0$
hasta $p_m$ es una secuencia
\[
  \pi = p_0 a_1 p_1 a_2 p_2 \cdots a_m p_m
\]
tal que, para cada $1 \le i \le m$, se cumple
\[
  \delta(p_{i-1},a_i)=p_i.
\]
La cadena que se lee a lo largo de $\pi$, denotada $w(\pi)$, es
$a_1a_2\cdots a_m$. Si $m=0$, el camino no contiene transiciones y la
cadena que se lee es $\varepsilon$.
\end{definicion}

\begin{definicion}
Un camino $\pi$ en el autómata es \textbf{aceptante} si empieza en el
estado inicial $q_0$ y termina en un estado de aceptación. El lenguaje
aceptado por $A$ es
\[
  L(A) =
  \{w(\pi) \mid \pi
  \text{ es un camino aceptante de } A\}.
\]
\end{definicion}

La aciclicidad se define a partir de los caminos del autómata.

\begin{definicion}
Un autómata finito $A$ es \textbf{acíclico} si
ningún camino de $A$ con al menos una transición empieza y termina en
el mismo estado.
\end{definicion}

La aciclicidad es suficiente para asegurar que el lenguaje aceptado es
finito. En los tres lemas siguientes se enuncian la cota de longitud de
los caminos, que se usa en el análisis de complejidad, la finitud del
lenguaje aceptado y la representación de todo lenguaje finito mediante
un autómata acíclico.

\begin{lema}
\label{lema:longitud-caminos-aciclicos}
Sea $A=(Q,\Sigma,\delta,q_0,F)$ un autómata finito acíclico. Todo
camino de $A$ tiene longitud a lo sumo $|Q|-1$. En
consecuencia, toda cadena aceptada por $A$ tiene longitud a lo sumo
$|Q|-1$~\cite{bertsch2001rcg}.
\end{lema}

\begin{lema}
\label{lema:lenguaje-finito-aciclico}
Si $A$ es un autómata finito acíclico, entonces
$L(A)$ es finito~\cite{bertsch2001rcg}.
\end{lema}

\begin{lema}
\label{lema:finito-reconocible-aciclico}
Todo lenguaje finito puede reconocerse mediante un autómata finito
acíclico~\cite{bertsch2001rcg}.
\end{lema}

En el análisis de complejidad del presente capítulo, el autómata forma
parte de la entrada, de modo que el costo se expresa en función de su
tamaño. Ese costo depende, ante todo, de la cantidad de caminos
aceptantes del autómata. Por ello se introduce la siguiente notación.

\begin{definicion}
Sea $A$ un autómata finito acíclico. Se denota por $p(A)$ la cantidad
de caminos aceptantes de $A$.
\end{definicion}

Como $\delta$ asigna a lo sumo un estado a cada par $(q,a)$, cada
cadena aceptada determina un único camino aceptante. En consecuencia,
\[
  p(A)=|L(A)|.
\]

\section{Problema de vacío de la intersección}
\label{sec:problema-vacio-interseccion}

Las secciones anteriores describieron la parte regular de la entrada. El
\textbf{problema del vacío de la intersección de una RCG con un lenguaje
regular finito} toma como entrada una RCG $G$ y un autómata
finito acíclico $A$, y pregunta si $L(G) \cap L(A) \neq \emptyset$.

La intersección es no vacía si y solo si existe una cadena que se
reconoce a la vez por la RCG $G$ y por el autómata $A$:
\[
  L(G) \cap L(A) \neq \emptyset
  \quad \Longleftrightarrow \quad
  \exists w \in \Sigma^* \text{ tal que } w \in L(G) \text{ y } w \in L(A).
\]

El algoritmo de la sección siguiente decide el problema recorriendo las
cadenas aceptadas por el autómata acíclico.

\section{Algoritmo de decisión}
\label{sec:algoritmo-decision}

El resultado de Bertsch y Nederhof garantiza que este caso es decidible,
pero no proporciona un algoritmo explícito ni expresa su costo en función
de la cantidad de cadenas aceptadas~\cite{bertsch2001rcg}. La presente
sección da ese algoritmo y analiza su costo.

Como el autómata de entrada es acíclico, sus caminos aceptantes son
finitos. Además, como $\delta$ asigna a lo sumo un estado a cada par
$(q,a)$, esos caminos están en correspondencia uno a uno con las
cadenas aceptadas. Por tanto, una
forma directa de decidir si la intersección es no vacía consiste en
recorrer las cadenas aceptadas por el autómata y aplicar, a cada una, un
algoritmo de reconocimiento para RCG. El procedimiento es
determinista y, en general, no se ejecuta en tiempo polinomial. Su utilidad está en
exponer el parámetro responsable del costo: la cantidad de cadenas
aceptadas por el autómata.

Se usará como subrutina un procedimiento de reconocimiento para RCG, que
se denota
\[
  \proc{RCG-Recognize}(G,w).
\]
La subrutina devuelve \const{true} si y solo si $w \in L(G)$. El
reconocimiento de una cadena por una RCG se resuelve en tiempo
polinomial en la longitud de la
cadena~\cite{bertsch2001rcg,kallmeyer2010parsing}.

Antes de enumerar cadenas conviene eliminar los estados que no pueden
participar en ningún camino aceptante. Esa reducción evita recorrer
ramas que no contribuyen al lenguaje del autómata.

\begin{definicion}
Sea $A=(Q,\Sigma,\delta,q_0,F)$ un autómata finito. Un estado $q\in Q$
es \textbf{útil} si es alcanzable desde $q_0$ y desde
$q$ se puede alcanzar algún estado de aceptación.
\end{definicion}

El conjunto de estados útiles se obtiene mediante dos recorridos del
autómata. En el primero se determinan los estados alcanzables desde
$q_0$: se parte de $\{q_0\}$ y se agrega $\delta(q,a)$ por cada estado
$q$ ya incluido y cada símbolo $a$ para el que $\delta(q,a)$ está
definida, hasta que no se agregue ningún estado. En el segundo se
determinan los estados desde los que se alcanza un estado de aceptación:
se parte de $F$ y se agrega cada estado $q$ tal que $\delta(q,a)$ ya
está incluido para algún símbolo $a$, hasta que no se agregue ninguno.
Un estado es útil cuando pertenece a los dos conjuntos. El autómata
$\proc{Useful-Part}(A)$ tiene por estados los útiles, por estado inicial
$q_0$, por estados de aceptación los de $F$ que son útiles y por
transiciones las $\delta(q,a)=q'$ en las que $q$ y $q'$ son útiles. En
cada recorrido, cada estado y cada transición se visitan a lo sumo una
vez, de modo que la reducción se calcula en tiempo lineal en el tamaño
del autómata.

El pseudocódigo de $\proc{Useful-Part}$ es el siguiente.

\begin{codebox}
\Procname{$\proc{Useful-Part}(A)$}
\li $R = \{q_0\}$
\li \While there is a transition $\delta(q,a)=q'$ with $q \in R$ and $q' \notin R$
\li \Do $R = R \cup \{q'\}$
    \End
\li $C = F$
\li \While there is a transition $\delta(q,a)=q'$ with $q' \in C$ and $q \notin C$
\li \Do $C = C \cup \{q\}$
    \End
\li $U = R \cap C$
\li \Return $(U, \Sigma, \delta', q_0, F \cap U)$, where $\delta'$ is the
    restriction of $\delta$ to $U$
\end{codebox}

El autómata $\proc{Useful-Part}(A)$ acepta el mismo lenguaje que $A$;
cuando no hay estados útiles, ese lenguaje es vacío.

El algoritmo principal se presenta a continuación.

\begin{codebox}
\Procname{$\proc{Nonempty-Intersection}(G,A)$}
\li $A' = \proc{Useful-Part}(A)$
\li \If the initial state of $A'$ does not exist
\li \Then \Return \const{false}
    \End
\li \Return $\proc{Search}(q_0,\varepsilon)$
\end{codebox}

La subrutina $\proc{Search}$ recorre las cadenas aceptadas por el
autómata. El parámetro $q$ indica el estado actual y $w$ almacena la
cadena que se lee a lo largo del camino desde $q_0$ hasta ese estado.

\begin{codebox}
\Procname{$\proc{Search}(q,w)$}
\li \If $q \in F$ \kw{and} $\proc{RCG-Recognize}(G,w)$
\li \Then \Return \const{true}
    \End
\li \For each symbol $a \in \Sigma$ such that $\delta'(q,a)$ is defined
\li \Do
        $r = \delta'(q,a)$
\li     \If $\proc{Search}(r,wa)$
\li     \Then \Return \const{true}
        \End
    \End
\li \Return \const{false}
\end{codebox}

En el pseudocódigo, $\delta'$ denota la función de transición de $A'$.
La línea~1 de $\proc{Search}$ verifica la pertenencia a la RCG cada vez
que el camino recorrido termina en un estado de aceptación. Si la
verificación
es positiva, la intersección es no vacía y el algoritmo termina. Si la
verificación es negativa, el recorrido continúa, ya que desde un estado
de aceptación pueden salir transiciones hacia otros caminos aceptantes.
La línea
5 actualiza la cadena asociada al camino mediante la concatenación del
símbolo leído.

La aciclicidad de $A$ garantiza la terminación de la búsqueda. Los
estados $q_0,q_1,\ldots,q_d$ que aparecen en una secuencia de llamadas
recursivas anidadas forman un camino de $A'$: en cada paso, de la
llamada en el estado $q$ se pasa a la llamada en el estado
$\delta'(q,a)$. La profundidad $d$ de la recursión es entonces el número
de transiciones de ese camino. Como $A'$ es acíclico, ese camino no
contiene dos veces el mismo estado, de modo que ningún estado se repite
a lo largo de la recursión. Por el
lema~\ref{lema:longitud-caminos-aciclicos}, $d \le |Q|-1$. Como en cada
llamada se generan a lo sumo $|\Sigma|$ llamadas recursivas, el árbol de
recursión es finito y la búsqueda termina.

\subsection{Correctitud}
\label{sec:correctitud-algoritmo}

La correctitud se demuestra separando tres hechos. El primero establece
que toda cadena que se evalúa con $\proc{RCG-Recognize}$ pertenece al
lenguaje del autómata. El segundo establece que toda cadena aceptada por
el autómata se evalúa exactamente una vez. El tercero se apoya en la
correctitud de $\proc{RCG-Recognize}$.

\begin{lema}
\label{lema:sonido-automata}
Toda cadena $w$ para la que el algoritmo ejecuta
$\proc{RCG-Recognize}(G,w)$ pertenece a $L(A)$.
\end{lema}

\begin{proof}
La subrutina $\proc{Search}$ ejecuta $\proc{RCG-Recognize}(G,w)$ solo en
la línea~1, cuando el estado actual $q$ pertenece a $F$. El valor de $w$
se construye concatenando los símbolos que se leen desde $q_0$ hasta
$q$. Por tanto, existe un camino desde $q_0$ hasta un estado de
aceptación a lo
largo del cual se lee $w$. Por la definición de aceptación,
$w \in L(A')$. Como $A'=\proc{Useful-Part}(A)$ acepta el mismo lenguaje
que $A$, se tiene $L(A')=L(A)$. Luego $w \in L(A)$.
\end{proof}

\begin{lema}
\label{lema:completo-automata}
Para toda cadena $w \in L(A)$, el algoritmo ejecuta
$\proc{RCG-Recognize}(G,w)$ exactamente una vez.
\end{lema}

\begin{proof}
Sea $w=a_1a_2\cdots a_m$ una cadena de $L(A)$. Como para cada par
$(q,a)$ hay a lo sumo un estado $\delta(q,a)$, existe un único camino de
$A$ que comienza en $q_0$ y a lo largo del cual se lee $w$:
\[
  q_0 a_1 q_1 a_2 q_2 \cdots a_m q_m.
\]
Además, como $w \in L(A)$, se tiene $q_m \in F$. Todos los estados de ese
camino son útiles: cada uno es alcanzable desde $q_0$ por el prefijo
correspondiente, y desde cada uno se alcanza un estado de aceptación
por el
sufijo correspondiente. Por tanto, el camino también pertenece a $A'$.

La subrutina $\proc{Search}$ recorre, para cada estado actual, todos los
símbolos $a$ para los cuales $\delta'(q,a)$ está definida. Por inducción
sobre la longitud de los prefijos de $w$, el algoritmo realiza una
llamada recursiva asociada al único prefijo de longitud $i$ de $w$, para
cada $0 \le i \le m$. Al llegar al estado $q_m$, la cadena acumulada es
$w$ y se cumple $q_m \in F$. En ese momento se ejecuta
$\proc{RCG-Recognize}(G,w)$.

La ejecución es única porque en $A'$ hay a lo sumo un estado
$\delta'(q,a)$ por cada par $(q,a)$: ninguna cadena se lee a lo largo
de dos caminos distintos desde $q_0$.
Por tanto, el algoritmo no puede alcanzar un estado de aceptación con
la misma
cadena $w$ por dos recorridos diferentes.
\end{proof}

\begin{lema}
\label{lema:correctitud-reconocedor}
Para toda cadena $w$, la subrutina $\proc{RCG-Recognize}(G,w)$ devuelve
\const{true} si y solo si $w \in L(G)$.
\end{lema}

\begin{proof}
El reconocimiento de una cadena por una RCG se decide en tiempo
polinomial~\cite{bertsch2001rcg,kallmeyer2010parsing}; $\proc{RCG-Recognize}$
es un procedimiento que resuelve ese problema y devuelve \const{true} si
y solo si $w \in L(G)$.
\end{proof}

\begin{teorema}
\label{teo:correctitud-enumeracion}
El algoritmo $\proc{Nonempty-Intersection}(G,A)$ devuelve
\const{true} si y solo si
\[
  L(G) \cap L(A) \neq \emptyset.
\]
\end{teorema}

\begin{proof}
Supóngase primero que el algoritmo devuelve \const{true}. Entonces, para
alguna cadena $w$, se ejecutó $\proc{RCG-Recognize}(G,w)$ y la subrutina
devolvió \const{true}. Por el lema~\ref{lema:sonido-automata},
$w \in L(A)$. Por el lema~\ref{lema:correctitud-reconocedor},
$w \in L(G)$. Luego $w \in L(G) \cap L(A)$, y la intersección es no
vacía.

Supóngase ahora que $L(G) \cap L(A) \neq \emptyset$. Entonces existe una
cadena $w$ tal que $w \in L(G)$ y $w \in L(A)$. Por el
lema~\ref{lema:completo-automata}, el algoritmo ejecuta
$\proc{RCG-Recognize}(G,w)$ exactamente una vez. Por el
lema~\ref{lema:correctitud-reconocedor}, esa ejecución devuelve
\const{true}. Por tanto, el algoritmo devuelve \const{true}.
\end{proof}

\subsection{Complejidad}
\label{sec:complejidad-algoritmo}

La complejidad del algoritmo depende de dos componentes. El primero es
la enumeración de cadenas aceptadas por el autómata acíclico. El
segundo es el reconocimiento de cada cadena mediante
la RCG. En la presente subsección se expresan ambos costos de manera
separada.

Sea $|A|=|Q|+|\delta|$ el tamaño del autómata, donde $|\delta|$ denota
la cantidad de pares $(q,a)$ para los cuales la transición está definida.
Sea $p(A)$ el número de caminos aceptantes de $A$. Como se estableció
en la sección~\ref{sec:automatas-aciclicos}, $p(A)=|L(A)|$.

Después de eliminar estados no útiles, todo camino recorrido por el
algoritmo es prefijo de algún camino aceptante. Como cada camino
aceptante tiene longitud a lo sumo $|Q|-1$, el número de prefijos
recorridos está acotado por
\[
  p(A)\,|Q|.
\]
Por tanto, el costo de la enumeración, sin contar las llamadas a
$\proc{RCG-Recognize}$, queda acotado por
\[
  O(p(A)\,|Q|),
\]
si se usan listas de adyacencia y la cadena acumulada se mantiene en
una pila de símbolos.

Sea $R_G(n)$ el tiempo que demora
$\proc{RCG-Recognize}(G,w)$ sobre una cadena $w$ de longitud $n$. Para
una RCG fija, $R_G(n)$ es polinomial en $n$. Si la gramática forma parte
de la entrada, el costo es polinomial en la longitud de la cadena para
cada gramática fijada, y el exponente depende de parámetros
estructurales de la gramática, como la cantidad de variables que pueden
recibir rangos en una cláusula~\cite{bertsch2001rcg,kallmeyer2010parsing}.

Por el lema~\ref{lema:longitud-caminos-aciclicos}, toda cadena aceptada
tiene longitud a lo sumo $|Q|-1$. Luego el costo total del algoritmo se
acota por
\[
  O\bigl( |A| + p(A)\,|Q| + p(A)\,R_G(|Q|-1) \bigr).
\]
De forma equivalente,
\[
  O\bigl( |A| + p(A)(|Q| + R_G(|Q|-1)) \bigr).
\]

La cota anterior es polinomial si $p(A)$ está acotado por un polinomio
en $|A|$. Sin embargo, $p(A)$ no está polinomialmente acotado para todos
los autómatas acíclicos. En la
sección~\ref{sec:lenguajes-regulares-finitos} ya se mostró un
autómata así: el que reconoce $L_n=\{0,1\}^n$, con $n+1$ estados,
acíclico y con $2^n$ cadenas aceptadas. Por la igualdad $p(A)=|L(A)|$,
ese autómata tiene $2^n$ caminos aceptantes con un tamaño lineal en $n$.

Por tanto, la enumeración de cadenas aceptadas puede tener
costo exponencial en el tamaño del autómata, incluso para autómatas
acíclicos. En consecuencia, el algoritmo del presente
capítulo no contradice la NP-completitud del caso finito con autómatas
acíclicos demostrada en~\cite{bertsch2001rcg}. La enumeración
determinista decide el problema, pero puede tardar tiempo exponencial.

La cota anterior permite obtener una condición suficiente de
tratabilidad para clases restringidas de autómatas acíclicos.

\begin{teorema}
\label{teo:clases-pol-caminos}
Sea $\mathcal{A}$ una clase de autómatas finitos acíclicos tal que
existe un polinomio $r$ con
\[
  p(A) \le r(|A|)
\]
para todo $A \in \mathcal{A}$. Entonces el problema del vacío de la
intersección, restringido a autómatas de $\mathcal{A}$, se decide en
tiempo polinomial para cada RCG fija.
\end{teorema}

\begin{proof}
Por la cota del algoritmo,
\[
  O\bigl( |A| + p(A)(|Q| + R_G(|Q|-1)) \bigr)
\]
es suficiente para decidir el problema. Si $A \in \mathcal{A}$,
entonces $p(A) \le r(|A|)$.
Además, $|Q| \le |A|$. Para una RCG fija, $R_G(|Q|-1)$ es polinomial en
$|Q|$, y por tanto en $|A|$. La composición de esos polinomios produce
un tiempo polinomial en $|A|$.
\end{proof}

El teorema anterior debe leerse como una condición sobre la
representación regular. Si una clase de autómatas acíclicos tiene
polinomialmente muchas cadenas aceptadas, la
enumeración alcanza para obtener una decisión polinomial. Si una clase
admite autómatas que representan exponencialmente muchas cadenas, el
algoritmo puede ser exponencial.

Generar cada cadena aceptada en tiempo polinomial no es lo mismo que
enumerarlas todas en tiempo polinomial. La cantidad de cadenas
aceptadas puede ser exponencial, de modo que recorrerlas todas cuesta
tiempo exponencial aun cuando cada una se obtenga en tiempo polinomial.
Para decidir el vacío de la intersección hace falta, en el peor caso,
recorrer todas: cuando la intersección es vacía, ninguna cadena aceptada
pertenece a $L(G)$, y el recorrido solo concluye al agotarlas.

\section{Modelación del SAT con RCG}
\label{sec:modelacion-sat}

El SAT se modela con el problema de vacío de las secciones anteriores. A
una fórmula con $n$ ocurrencias de variable se le asocia el autómata
booleano del capítulo~\ref{cap:rcg-sat}, finito y acíclico, que reconoce
sus interpretaciones verdaderas bajo el supuesto de que cada ocurrencia
es independiente. Queda imponer que las ocurrencias de una misma
variable tomen el mismo valor. Tras SATCF y SATSRC, este
trabajo analiza esa consistencia con una RCG.

Sea $\ell_1\cdots\ell_n$ la secuencia de ocurrencias de la fórmula, con
$\ell_i$ la variable de la ocurrencia $i$. La gramática de consistencia
$G_{\mathrm{cons}}$ tiene por terminales $T=\{0,1\}$, por variables
$X, W_0, W_1, \ldots, W_n$ y tres juegos de predicados, todos de
aridad uno: el predicado inicial $S$, un predicado $Eq_v$ por cada
variable booleana $v$ y los predicados $L_0, L_1, \ldots, L_n$. El
predicado $L_j$ reconoce las cadenas de longitud $j$, mediante las
cláusulas
\[
  L_j(cX) \to L_{j-1}(X) \quad (c \in \{0,1\},\ 1 \le j \le n),
  \qquad L_0(\varepsilon) \to \varepsilon.
\]

Para cada variable booleana $v$, con ocurrencias en las posiciones
$i_1 < i_2 < \cdots < i_r$, el predicado $Eq_v$ reconoce las cadenas de
longitud $n$ en las que esas posiciones llevan el mismo símbolo. Sus
cláusulas son dos, una por cada símbolo $b \in \{0,1\}$:
\[
  Eq_v(W_0\, b\, W_1\, b \cdots b\, W_r) \to
  L_{g_0}(W_0)\, L_{g_1}(W_1) \cdots L_{g_r}(W_r).
\]
El argumento es un patrón: el símbolo $b$ ocupa las posiciones de $v$, y
cada variable $W_k$ cubre una subcadena: la anterior a la primera
ocurrencia, la comprendida entre dos ocurrencias consecutivas o la
posterior a la última. Esas subcadenas tienen longitudes
$g_0 = i_1 - 1$, $g_k = i_{k+1} - i_k - 1$ y $g_r = n - i_r$, y cada
una se reconoce por el predicado $L_{g_k}$ correspondiente. Una vez
fijada la fórmula, las posiciones $i_k$ y las longitudes $g_k$ son
constantes, de modo que el conjunto de cláusulas es finito y cada
cláusula tiene la forma de la definición~\ref{def:rcg}.

La consistencia
de la fórmula es la conjunción de esas restricciones, y se expresa con
una cláusula no lineal con el mismo argumento en los $m$ predicados
$Eq_v$, donde $m$ es la cantidad de variables booleanas de la fórmula:
\[
  S(X) \to Eq_{v_1}(X)\, Eq_{v_2}(X) \cdots Eq_{v_m}(X).
\]
La variable $X$ es la misma en los $m$ predicados, así que las $m$
restricciones se aplican a la misma cadena. Con esta cláusula queda
completo el conjunto de cláusulas y, con él, la gramática
$G_{\mathrm{cons}}$. El lenguaje reconocido por $G_{\mathrm{cons}}$ se
denota $L(G_{\mathrm{cons}})$; el lema siguiente lo caracteriza.

\begin{lema}
\label{lema:lenguaje-gcons}
Sea $G_{\mathrm{cons}}$ la gramática de consistencia de una fórmula
booleana con $n$ ocurrencias de variable. El lenguaje
$L(G_{\mathrm{cons}})$ es el conjunto de las cadenas consistentes de
la fórmula.
\end{lema}

\begin{proof}
Primero se establece que el predicado $L_j$ deriva en la cadena vacía
sobre una cadena $u$ si y solo si $|u| = j$. La prueba es por
inducción sobre $j$: la cláusula $L_0(\varepsilon) \to \varepsilon$
cubre el caso $j = 0$, y con cada cláusula $L_j(cX) \to L_{j-1}(X)$
se consume un símbolo.

Sea $w$ una cadena consistente; como lleva un símbolo por ocurrencia,
$|w| = n$. Para cada variable booleana $v$, sea $b$ el símbolo común
de sus posiciones en $w$. Entonces $w$ se descompone como
$u_0\, b\, u_1\, b \cdots b\, u_r$ con $|u_k| = g_k$: en el patrón
del símbolo $b$, la sustitución de rango asocia cada variable $W_k$ a
la subcadena $u_k$, y cada predicado $L_{g_k}(u_k)$ deriva en la
cadena vacía. Luego $Eq_v(w)$ deriva en la cadena vacía. Como esto
vale para las $m$ variables booleanas, con la cláusula de $S$ se
deriva en la cadena vacía y $w \in L(G_{\mathrm{cons}})$.

Sea ahora $w \in L(G_{\mathrm{cons}})$. La única cláusula de $S$ es
$S(X) \to Eq_{v_1}(X) \cdots Eq_{v_m}(X)$, de modo que cada $Eq_v(w)$
deriva en la cadena vacía. En esa derivación se aplica una de las dos
cláusulas de $Eq_v$, con algún símbolo $b$: entonces
$w = u_0\, b\, u_1\, b \cdots b\, u_r$, y cada predicado
$L_{g_k}(u_k)$ deriva en la cadena vacía, así que $|u_k| = g_k$. Como
$g_0 + g_1 + \cdots + g_r + r = n$ por la definición de las
longitudes $g_k$, se tiene $|w| = n$, y las posiciones
$i_1, \ldots, i_r$ de $w$ llevan el símbolo $b$. Como esto vale para
cada variable booleana, $w$ es consistente.
\end{proof}

Por ejemplo, considere la fórmula $x_1 \wedge (x_2 \vee x_1)$ del
capítulo~\ref{cap:rcg-sat}, cuya secuencia de ocurrencias es
$x_1\,x_2\,x_1$, con $n=3$ y $m=2$. La variable $x_1$ ocurre en las
posiciones $1$ y $3$, con subcadenas de longitudes $g_0=0$, $g_1=1$ y
$g_2=0$; la variable $x_2$ ocurre en la posición $2$, con subcadenas
de longitudes $g_0=1$ y $g_1=1$. Las cláusulas de $G_{\mathrm{cons}}$
son las siguientes, omitidas las de $L_2$ y $L_3$, que en ningún
patrón aparecen:
\[
\begin{aligned}
  S(X) &\to Eq_{x_1}(X)\, Eq_{x_2}(X),\\
  Eq_{x_1}(W_0\,0\,W_1\,0\,W_2) &\to L_0(W_0)\, L_1(W_1)\, L_0(W_2),\\
  Eq_{x_1}(W_0\,1\,W_1\,1\,W_2) &\to L_0(W_0)\, L_1(W_1)\, L_0(W_2),\\
  Eq_{x_2}(W_0\,0\,W_1) &\to L_1(W_0)\, L_1(W_1),\\
  Eq_{x_2}(W_0\,1\,W_1) &\to L_1(W_0)\, L_1(W_1),\\
  L_1(0X) \to L_0(X), \quad L_1(1X) &\to L_0(X), \qquad
  L_0(\varepsilon) \to \varepsilon.
\end{aligned}
\]

La cadena $101$ es consistente y se reconoce por $G_{\mathrm{cons}}$
mediante las derivaciones siguientes:
\[
\begin{aligned}
  S(101) &\to Eq_{x_1}(101)\, Eq_{x_2}(101),\\
  Eq_{x_1}(101) &\to L_0(\varepsilon)\, L_1(0)\, L_0(\varepsilon),\\
  Eq_{x_2}(101) &\to L_1(1)\, L_1(1).
\end{aligned}
\]
A continuación se describen los pasos.
\begin{itemize}
  \item \textbf{Primer paso:} en la cláusula
    $S(X) \to Eq_{x_1}(X)\, Eq_{x_2}(X)$, la sustitución de rango
    asocia la variable $X$ al valor $X = 101$. De esta forma se deriva
    en $Eq_{x_1}(101)\, Eq_{x_2}(101)$.
  \item \textbf{Segundo paso:} en el patrón del símbolo $1$ de
    $Eq_{x_1}$, la sustitución de rango asocia las variables $W_0$,
    $W_1$, $W_2$ a los valores $W_0 = \varepsilon$, $W_1 = 0$ y
    $W_2 = \varepsilon$. Con estos valores se deriva en
    $L_0(\varepsilon)\, L_1(0)\, L_0(\varepsilon)$.
  \item \textbf{Tercer paso:} en el patrón del símbolo $0$ de
    $Eq_{x_2}$, la sustitución de rango asocia las variables $W_0$,
    $W_1$ a los valores $W_0 = 1$ y $W_1 = 1$. Con estos valores se
    deriva en $L_1(1)\, L_1(1)$.
  \item \textbf{Cuarto paso:} cada predicado $L_1$ deriva en
    $L_0(\varepsilon)$, y cada predicado $L_0(\varepsilon)$ deriva en
    la cadena vacía por la cláusula terminal, con lo que la cadena
    $101$ queda reconocida y $101 \in L(G_{\mathrm{cons}})$.
\end{itemize}

La cadena $110$, con la que en el capítulo~\ref{cap:rcg-sat} se llega
al estado $V$ del autómata booleano, no es consistente. En
$Eq_{x_1}(110)$ ningún patrón admite una sustitución de rango: el del
símbolo $0$ exige $0$ en la posición $1$, donde la cadena lleva $1$, y
el del símbolo $1$ exige $1$ en la posición $3$, donde la cadena lleva
$0$.
Como $Eq_{x_1}(110)$ no deriva en la cadena vacía,
$110 \notin L(G_{\mathrm{cons}})$.

La satisfacibilidad de la fórmula equivale a que la intersección de
ambos lenguajes no sea vacía, como establece el teorema siguiente.

\begin{teorema}
\label{teo:sat-vacio-interseccion}
Sean $G_{\mathrm{cons}}$ la gramática de consistencia de una fórmula
booleana y $A$ su autómata booleano. La fórmula es satisfacible si y
solo si $L(G_{\mathrm{cons}}) \cap L(A) \neq \emptyset$.
\end{teorema}

\begin{proof}
El lenguaje del autómata booleano está formado por las cadenas que
hacen verdadera la fórmula cuando cada ocurrencia se trata por
separado~\cite{fernandezarias2007}.

A cada interpretación $I$ le corresponde la cadena
$w_I = b_1 b_2 \cdots b_n$, donde $b_i$ es el valor, $1$ o $0$, que
$I$ asigna a la variable $\ell_i$. La cadena $w_I$ es consistente,
porque las posiciones de una misma variable llevan el valor que $I$
asigna a esa variable, y toda cadena consistente es $w_I$ para la
interpretación $I$ que asigna a cada variable el símbolo común de sus
posiciones. Además, en $w_I$ cada ocurrencia lleva el valor que $I$
asigna a su variable, de modo que $w_I \in L(A)$ si y solo si $I$
hace verdadera la fórmula.

Supóngase que la fórmula es satisfacible. Entonces alguna
interpretación $I$ la hace verdadera. La cadena $w_I$ es consistente,
de modo que $w_I \in L(G_{\mathrm{cons}})$ por el
lema~\ref{lema:lenguaje-gcons}, y $w_I \in L(A)$. Luego la
intersección es no vacía.

Supóngase ahora que la intersección es no vacía. Entonces alguna
cadena $w$ pertenece a $L(G_{\mathrm{cons}})$ y a $L(A)$. Por el
lema~\ref{lema:lenguaje-gcons}, $w$ es consistente, de modo que
$w = w_I$ para la interpretación $I$ que asigna a cada variable el
símbolo común de sus posiciones. De $w_I \in L(A)$ se sigue que $I$
hace verdadera la fórmula. Por tanto, la fórmula es satisfacible.
\end{proof}

El SAT se reduce entonces al problema de vacío de la intersección que
este capítulo resuelve.

\section{Consecuencias para el SAT y subclases tratables}
\label{sec:lectura-sat}

Con la enumeración de este capítulo se decide esa satisfacibilidad: se
recorren las cadenas aceptadas por el autómata booleano y se aplica a
cada una el reconocimiento de la RCG de consistencia. Las dos partes
de la instancia tienen costo polinomial. El autómata booleano tiene
$O(n^2)$ estados~\cite{fernandezarias2007}, con $n$ la cantidad de
ocurrencias de variables booleanas en la fórmula. La gramática
$G_{\mathrm{cons}}$ tiene aridad uno y tamaño $O(nm)$, contado en
símbolos, con $m$ la cantidad de variables booleanas distintas: dos
cláusulas de longitud $O(n)$ por variable booleana, más las $O(n)$
cláusulas de los predicados $L_j$. El reconocimiento con
$G_{\mathrm{cons}}$ es polinomial: cada cláusula admite a lo sumo una sustitución de rango
para cada cadena, porque los rangos de las variables $W_k$ quedan
determinados por las longitudes $g_k$, y esa sustitución se comprueba
en tiempo lineal; el reconocimiento de una cadena de longitud $n$
cuesta entonces $R_{G_{\mathrm{cons}}}(n) = O(nm)$.

Toda cadena aceptada por el autómata booleano tiene longitud $n$, de
modo que cada llamada a $\proc{RCG-Recognize}$ cuesta $O(nm)$. Con
$|Q| = O(n^2)$, la cota de la
subsección~\ref{sec:complejidad-algoritmo} queda en
\[
  O\bigl( n^2 + p(A)\,(n^2 + nm) \bigr) = O\bigl( n^2 + p(A)\,n^2 \bigr),
\]
pues $m \le n$. El factor dominante es $p(A)$, la cantidad de cadenas
aceptadas por el autómata booleano, que para una fórmula con muchas
interpretaciones verdaderas bajo el supuesto de independencia es del
orden de $2^n$. La dificultad no está entonces en el tamaño de la
instancia ni en el reconocimiento, sino en el vacío: decidir si
$L(G_{\mathrm{cons}}) \cap L(A)$ es vacío es el SAT, y para una RCG
con un autómata finito acíclico ese problema es
NP-completo~\cite{bertsch2001rcg}.

En SATCF y SATSRC ese costo no aparece, y la razón no es una enumeración
más hábil. En esos enfoques el vacío se decide desde el formalismo: el
de una gramática libre del contexto y el de una sRCG se deciden en
tiempo polinomial, sin recorrer las
interpretaciones~\cite{fernandezarias2007,aguileralopez2016}. Para las
RCG no hay ese procedimiento directo: con la enumeración el problema se
decide, pero el costo es el número de interpretaciones, no una propiedad
del formalismo.

La no-linealidad tiene dos efectos. Por un lado, con la igualdad
entre las subcadenas en que aparece una misma variable se amplía la
clase de fórmulas cuya consistencia se expresa, más allá del orden de
concatenación de rango simple de las sRCG. Por otro, por esa misma
igualdad el vacío de la intersección deja de ser polinomial.

Quedan entonces dos vías hacia una extensión tratable de SATSRC. En
la primera se restringen las fórmulas: si el autómata booleano tiene
una cantidad de caminos aceptantes acotada por un polinomio en $n$,
la cota anterior es polinomial y el algoritmo de enumeración
decide la satisfacibilidad en tiempo polinomial, por lo que ya esa
propuesta queda cerrada en este capítulo.

En la segunda vía se restringe la gramática de consistencia a una
subclase no lineal de las RCG, que debe conservar parte de la
expresividad que la no-linealidad añade sobre las sRCG y admitir un
procedimiento polinomial para el vacío de la intersección que no
recorra las interpretaciones. Con esta vía también se resuelve solo
una clase de fórmulas, aquellas cuyas cadenas consistentes se
reconocen por una gramática de la subclase: igual que en SATCF y
SATSRC, la clase de fórmulas queda delimitada por el formalismo de la
gramática de consistencia. Para hallar esa subclase hace falta una
representación simbólica de la intersección, sobre la que estudiar
con qué restricciones el vacío se vuelve tratable.

En el capítulo siguiente se construye esa representación: una gramática
producto entre una RCG y un lenguaje regular. Con ella el problema
general no se vuelve polinomial, pero se obtiene el objeto sobre el que
estudiar las subclases no lineales con vacío polinomial.


\bibliographystyle{plain}
\bibliography{../bib/referencias}

\end{document}
